\documentclass[UTF8,12pt,a4paper]{ctexart}
\usepackage[a4paper,margin=2cm]{geometry}
\usepackage{amsmath,amssymb,bm,booktabs,array,longtable,mathtools}
\usepackage{enumitem}
\usepackage{indentfirst}
\usepackage{microtype}
\linespread{1.4}
\setlength{\parindent}{2em}
\setlength{\parskip}{0em}
\setlist[itemize]{leftmargin=2em,itemsep=0.02em,topsep=0.04em,parsep=0em,partopsep=0em}
\setlist[enumerate]{leftmargin=2.2em,itemsep=0.03em,topsep=0.04em,parsep=0em,partopsep=0em}
\renewcommand\arraystretch{1.04}
\allowdisplaybreaks
\sloppy
\raggedbottom
\setlength{\jot}{0.08em}
\AtBeginDocument{\setlength{\abovedisplayskip}{0.16em}\setlength{\belowdisplayskip}{0.16em}\setlength{\abovedisplayshortskip}{0.08em}\setlength{\belowdisplayshortskip}{0.08em}}
\ctexset{section={format=\large\bfseries,beforeskip=0.45ex,afterskip=0.22ex},subsection={format=\normalsize\bfseries,beforeskip=0.35ex,afterskip=0.12ex}}
\newcommand{\dd}{\mathrm{d}}
\newcommand{\E}{\mathrm{E}}
\newcommand{\Var}{\mathrm{Var}}
\newcommand{\Cov}{\mathrm{Cov}}
\newcommand{\T}{\mathrm{T}}
\newcommand{\diag}{\mathrm{diag}}
\newcommand{\abs}[1]{\left|#1\right|}
\begin{document}

\begin{center}
    {\huge \textbf{误差理论与数据处理}}
\end{center}
\vspace{1em}

\section{随机误差与残余误差公式推导}

\subsection{测量误差、残余误差与平均值}
设被测量真值为 $X$，第 $i$ 次测量值为 $l_i$，随机误差为 $\delta_i=l_i-X$。若共有 $n$ 次等精度测量，则算术平均值为 $\bar{x}=\dfrac{1}{n}\sum_{i=1}^{n}l_i$。残余误差通常取为测量值与算术平均值之差，即 $v_i=l_i-\bar{x}$。于是有
\begin{equation*}
    \sum_{i=1}^{n}v_i=\sum_{i=1}^{n}(l_i-\bar{x})
    =\sum_{i=1}^{n}l_i-n\bar{x} .
\end{equation*}
若 $\bar{x}$ 由原始数据精确算出，满足 $n\bar{x}=\sum l_i$，所以
\begin{equation*}
    \sum_{i=1}^{n}v_i=0 .
\end{equation*}
这就是残余误差代数和为零的来源。

\subsection{残余误差代数和的凑整校核}
实际计算时，$\bar{x}$ 常常需要凑整。设精确平均值为 $\bar{x}_0$，凑整后的平均值为 $\bar{x}$，二者相差
\begin{equation*}
    \varepsilon=\bar{x}-\bar{x}_0,\qquad
    \sum_{i=1}^{n}v_i
    =\sum_{i=1}^{n}(l_i-\bar{x})
    =\sum_{i=1}^{n}l_i-n\bar{x}
    =n\bar{x}_0-n\bar{x}
    =-n\varepsilon .
\end{equation*}
若平均值按末位单位 $A$ 凑整，则 $\abs{\varepsilon}\le A/2$。因此 $\abs{\sum_{i=1}^{n}v_i}=n\abs{\varepsilon}\le nA/2$。若 $n$ 为偶数，$\sum v_i$ 的最大可能绝对值可达到 $nA/2$，故 $\abs{\sum_{i=1}^{n}v_i}\le nA/2$；若 $n$ 为奇数，$\sum v_i$ 通常只能取末位单位 $A$ 的整数倍，而 $nA/2$ 是半整数倍的 $A$，所以最大不超过 $(n-1)A/2$，即
\begin{equation*}
    \abs{\sum_{i=1}^{n}v_i}\le \frac{n-1}{2}A=\left(\frac{n}{2}-0.5\right)A .
\end{equation*}
这就是偶数、奇数两个校核公式的来源。

\subsection{真误差标准差公式}
若真值 $X$ 已知，随机误差为 $\delta_i=l_i-X$。等精度测量中，随机误差的数学期望为零，即
\begin{equation*}
    \E(\delta)=0 .
\end{equation*}
标准差由方差开方得到：
\begin{equation*}
    \sigma=\sqrt{\Var(\delta)}=\sqrt{\E(\delta^2)-[\E(\delta)]^2}.
\end{equation*}
由于 $\E(\delta)=0$，所以
\begin{equation*}
    \sigma=\sqrt{\E(\delta^2)} .
\end{equation*}
用有限次测量的平均平方误差代替数学期望，有
\begin{equation*}
    \sigma=\sqrt{\frac{1}{n}\sum_{i=1}^{n}\delta_i^2} .
\end{equation*}
该式适用于真值已知或可认为已知的情形。

\newpage

\subsection{贝塞尔公式推导}
真值未知时，不能直接计算 $\delta_i=l_i-X$，只能使用残余误差
\begin{equation*}
    v_i=l_i-\bar{x} .
\end{equation*}
由
\begin{equation*}
    l_i-X=(l_i-\bar{x})+(\bar{x}-X)=v_i+(\bar{x}-X)
\end{equation*}
得
\begin{equation*}
\begin{aligned}
    \sum_{i=1}^{n}(l_i-X)^2
    &=\sum_{i=1}^{n}\left[v_i+(\bar{x}-X)\right]^2 \\
    &=\sum_{i=1}^{n}v_i^2
      +2(\bar{x}-X)\sum_{i=1}^{n}v_i
      +n(\bar{x}-X)^2 .
\end{aligned}
\end{equation*}
由于 $\sum v_i=0$，中间项消失，因此
\begin{equation*}
    \sum_{i=1}^{n}(l_i-X)^2=\sum_{i=1}^{n}v_i^2+n(\bar{x}-X)^2 .
\end{equation*}
对两边取数学期望。因为 $\E[(l_i-X)^2]=\sigma^2$，所以
\begin{equation*}
    \E\left[\sum_{i=1}^{n}(l_i-X)^2\right]=n\sigma^2 .
\end{equation*}
又因为
\begin{equation*}
    \Var(\bar{x})=\frac{\sigma^2}{n},
\end{equation*}
所以
\begin{equation*}
    \E\left[n(\bar{x}-X)^2\right]=n\Var(\bar{x})=\sigma^2 .
\end{equation*}
因此
\begin{equation*}
    n\sigma^2=\E\left(\sum_{i=1}^{n}v_i^2\right)+\sigma^2,
\end{equation*}
即
\begin{equation*}
    \E\left(\sum_{i=1}^{n}v_i^2\right)=(n-1)\sigma^2 .
\end{equation*}
于是得到无偏估计式
\begin{equation*}
    s^2=\frac{\sum_{i=1}^{n}v_i^2}{n-1},
    \qquad
    s=\sqrt{\frac{\sum_{i=1}^{n}v_i^2}{n-1}} .
\end{equation*}
这就是贝塞尔公式。分母不是 $n$，而是 $n-1$，原因是平均值 $\bar{x}$ 已由数据估计出来，残余误差满足 $\sum v_i=0$，只有 $n-1$ 个独立信息量。

\newpage

\subsection{算术平均值标准差推导}
设 $l_1,l_2,\ldots,l_n$ 相互独立，且方差均为 $\sigma^2$。算术平均值为
\begin{equation*}
    \bar{x}=\frac{1}{n}\sum_{i=1}^{n}l_i .
\end{equation*}
根据方差性质，常数可以平方提出，独立变量和的方差等于方差之和，所以
\begin{equation*}
    \Var(\bar{x})=\Var\left(\frac{1}{n}\sum_{i=1}^{n}l_i\right)=\frac{1}{n^2}\sum_{i=1}^{n}\Var(l_i)=\frac{\sigma^2}{n}.
\end{equation*}
故算术平均值标准差为
\begin{equation*}
    \sigma_{\bar{x}}=\sqrt{\Var(\bar{x})}=\frac{\sigma}{\sqrt{n}} .
\end{equation*}
用贝塞尔公式估计 $\sigma$，可得
\begin{equation*}
    \sigma_{\bar{x}}=\sqrt{\frac{\sum_{i=1}^{n}v_i^2}{n(n-1)}} .
\end{equation*}

\subsection{或然误差与平均误差系数来源}
若随机误差服从正态分布 $N(0,\sigma^2)$，则误差落在 $[-\rho,\rho]$ 内的概率为 $1/2$ 时，$\rho$ 称为或然误差。令 $Z=\delta/\sigma$，则 $P(-\rho\le \delta\le \rho)=P(-\rho/\sigma\le Z\le \rho/\sigma)=1/2$。查标准正态分布表得 $\rho/\sigma=0.6745$，工程计算中常取 $0.6745\approx 2/3$，于是
\begin{equation*}
    \rho\approx \frac{2}{3}\sigma
    =\frac{2}{3}\sqrt{\frac{\sum_{i=1}^{n}v_i^2}{n-1}} .
\end{equation*}
平均误差常指绝对误差的数学期望。若 $\delta\sim N(0,\sigma^2)$，则
\begin{equation*}
\begin{aligned}
    \theta
    &=\E\abs{\delta}
      =\int_{-\infty}^{+\infty}\abs{x}\frac{1}{\sigma\sqrt{2\pi}}
        e^{-x^2/(2\sigma^2)}\dd x       \\
    &=2\int_{0}^{+\infty}x\frac{1}{\sigma\sqrt{2\pi}}
        e^{-x^2/(2\sigma^2)}\dd x       \\
    &=\sigma\sqrt{\frac{2}{\pi}}
      \approx 0.7979\sigma
      \approx \frac{4}{5}\sigma .
\end{aligned}
\end{equation*}
故有
\begin{equation*}
    \theta\approx \frac{4}{5}\sqrt{\frac{\sum_{i=1}^{n}v_i^2}{n-1}} .
\end{equation*}
算术平均值的或然误差与平均误差分别为
\begin{equation*}
    R=\frac{\rho}{\sqrt{n}},
    \qquad
    T=\frac{\theta}{\sqrt{n}} .
\end{equation*}
代入上式，得到
\begin{equation*}
    R=\frac{2}{3}\sqrt{\frac{\sum v_i^2}{n(n-1)}},
    \qquad
    T=\frac{4}{5}\sqrt{\frac{\sum v_i^2}{n(n-1)}} .
\end{equation*}

\newpage

\subsection{别捷尔斯法公式推导}
正态分布下，真误差的平均绝对值满足
\begin{equation*}
    \E\abs{\delta}=\sigma\sqrt{\frac{2}{\pi}} .
\end{equation*}
若使用残余误差 $v_i=l_i-\bar{x}$，则 $v_i$ 的方差不是 $\sigma^2$，而是
\begin{equation*}
    \Var(v_i)=\Var(l_i-\bar{x})=\sigma^2\left(1-\frac{1}{n}\right)=\frac{n-1}{n}\sigma^2 .
\end{equation*}
因此可近似认为
\begin{equation*}
    \E\abs{v_i}\approx \sqrt{\frac{n-1}{n}}\sigma\sqrt{\frac{2}{\pi}} .
\end{equation*}
将 $\E\abs{v_i}$ 用样本平均值 $\dfrac{1}{n}\sum\abs{v_i}$ 代替，得
\begin{equation*}
    \frac{1}{n}\sum_{i=1}^{n}\abs{v_i}
    \approx \sqrt{\frac{n-1}{n}}\sigma\sqrt{\frac{2}{\pi}} .
\end{equation*}
解出 $\sigma$：
\begin{equation*}
    \sigma\approx\sqrt{\frac{\pi}{2}}\sqrt{\frac{n}{n-1}}\frac{1}{n}\sum_{i=1}^{n}\abs{v_i}=\sqrt{\frac{\pi}{2}}\frac{\sum_{i=1}^{n}\abs{v_i}}{\sqrt{n(n-1)}}.
\end{equation*}
由于 $\sqrt{\pi/2}=1.253\cdots$，所以
\begin{equation*}
    \sigma=1.253\frac{\sum_{i=1}^{n}\abs{v_i}}{\sqrt{n(n-1)}} .
\end{equation*}

\section{不等精度测量公式推导}

\subsection{权与方差倒数成正比的推导}
设有 $m$ 个独立测量结果 $x_1,x_2,\ldots,x_m$，它们都测量同一个真值 $X$，且
\begin{equation*}
    \E(x_i)=X,
    \qquad
    \Var(x_i)=\sigma_i^2 .
\end{equation*}
构造线性无偏估计
\begin{equation*}
    \hat{X}=\sum_{i=1}^{m}a_i x_i .
\end{equation*}
无偏条件为
\begin{equation*}
    \E(\hat{X})=\sum_{i=1}^{m}a_i\E(x_i)=X\sum_{i=1}^{m}a_i=X,
\end{equation*}
所以必须满足
\begin{equation*}
    \sum_{i=1}^{m}a_i=1 .
\end{equation*}
由于各测量结果独立，估计量方差为
\begin{equation*}
    \Var(\hat{X})=\sum_{i=1}^{m}a_i^2\sigma_i^2 .
\end{equation*}
为了让估计最可靠，应在约束 $\sum a_i=1$ 下使方差最小。构造拉格朗日函数
\begin{equation*}
    \Phi=\sum_{i=1}^{m}a_i^2\sigma_i^2-\lambda\left(\sum_{i=1}^{m}a_i-1\right) .
\end{equation*}
对 $a_i$ 求偏导并令其为零：
\begin{equation*}
    \frac{\partial\Phi}{\partial a_i}=2a_i\sigma_i^2-\lambda=0 .
\end{equation*}
所以
\begin{equation*}
    a_i=\frac{\lambda}{2\sigma_i^2} .
\end{equation*}
代入 $\sum a_i=1$，得
\begin{equation*}
    \frac{\lambda}{2}\sum_{i=1}^{m}\frac{1}{\sigma_i^2}=1,
    \qquad
    \frac{\lambda}{2}=\frac{1}{\sum_{i=1}^{m}1/\sigma_i^2} .
\end{equation*}
于是
\begin{equation*}
    a_i=\frac{1/\sigma_i^2}{\sum_{j=1}^{m}1/\sigma_j^2} .
\end{equation*}
这说明权应与方差成反比，即
\begin{equation*}
    p_i\propto \frac{1}{\sigma_i^2} .
\end{equation*}
如果 $x_i$ 本身是第 $i$ 组测量的平均值，则
\begin{equation*}
    p_i\propto \frac{1}{\sigma_{\bar{x}_i}^2},
\end{equation*}
即
\begin{equation*}
    p_1:p_2:\cdots:p_m
    =\frac{1}{\sigma_{\bar{x}_1}^2}:\frac{1}{\sigma_{\bar{x}_2}^2}:\cdots:\frac{1}{\sigma_{\bar{x}_m}^2} .
\end{equation*}

\subsection{加权算术平均值推导}
由上节得到
\begin{equation*}
    \hat{X}=\sum_{i=1}^{m}a_i x_i,
    \qquad
    a_i=\frac{p_i}{\sum_{j=1}^{m}p_j} .
\end{equation*}
因此
\begin{equation*}
    \bar{x}=\hat{X}=\frac{\sum_{i=1}^{m}p_i x_i}{\sum_{i=1}^{m}p_i} .
\end{equation*}
为了减小数字运算量，常引入接近各 $x_i$ 的辅助值 $x_0$。因为
\begin{equation*}
    x_i=x_0+(x_i-x_0),
\end{equation*}
所以
\begin{equation*}
\begin{aligned}
    \bar{x}
    &=\frac{\sum_{i=1}^{m}p_i\left[x_0+(x_i-x_0)\right]}{\sum_{i=1}^{m}p_i} \\
    &=\frac{x_0\sum_{i=1}^{m}p_i+\sum_{i=1}^{m}p_i(x_i-x_0)}{\sum_{i=1}^{m}p_i} \\
    &=x_0+\frac{\sum_{i=1}^{m}p_i(x_i-x_0)}{\sum_{i=1}^{m}p_i} .
\end{aligned}
\end{equation*}
这就是加权平均值的常用计算形式。

\newpage

\subsection{加权平均值标准差推导}
若权的定义采用
\begin{equation*}
    p_i=\frac{\sigma_0^2}{\sigma_i^2},
\end{equation*}
其中 $\sigma_0$ 为单位权标准差，则
\begin{equation*}
    \sigma_i^2=\frac{\sigma_0^2}{p_i} .
\end{equation*}
加权平均值为
\begin{equation*}
    \bar{x}=\sum_{i=1}^{m}\frac{p_i}{\sum p_i}x_i .
\end{equation*}
由于各 $x_i$ 独立，故
\begin{equation*}
    \Var(\bar{x})=\sum_{i=1}^{m}\left(\frac{p_i}{\sum p_i}\right)^2\Var(x_i)=\frac{\sigma_0^2}{(\sum p_i)^2}\sum_{i=1}^{m}p_i=\frac{\sigma_0^2}{\sum_{i=1}^{m}p_i}.
\end{equation*}
所以
\begin{equation*}
    \sigma_{\bar{x}}=\frac{\sigma_0}{\sqrt{\sum_{i=1}^{m}p_i}} .
\end{equation*}
若 $\sigma_0$ 未知，用加权残差估计单位权方差：
\begin{equation*}
    s_0^2=\frac{\sum_{i=1}^{m}p_i v_i^2}{m-1},
    \qquad
    v_i=x_i-\bar{x} .
\end{equation*}
这是因为只估计了一个未知量 $X$，自由度为 $m-1$。代入得
\begin{equation*}
    \sigma_{\bar{x}}
    =\frac{s_0}{\sqrt{\sum p_i}}
    =\sqrt{\frac{\sum_{i=1}^{m}p_i v_i^2}{(m-1)\sum_{i=1}^{m}p_i}} .
\end{equation*}

\section{常见误差分布标准差推导}

\subsection{统一思路}
若误差 $\delta$ 的概率密度为 $f(\delta)$，则均值和方差为
\begin{equation*}
    \mu=\E(\delta)=\int \delta f(\delta)\dd\delta,
    \qquad
    \sigma^2=\Var(\delta)=\int(\delta-\mu)^2f(\delta)\dd\delta .
\end{equation*}
若分布关于零点对称，则 $\mu=0$，此时
\begin{equation*}
    \sigma^2=\int \delta^2 f(\delta)\dd\delta .
\end{equation*}

\subsection{均匀分布}
均匀分布为
\begin{equation*}
f(\delta)=\frac{1}{2a},\qquad -a\le \delta\le a .
\end{equation*}
该分布关于零点对称，故 $\mu=0$。于是
\begin{equation*}
\sigma^2=\int_{-a}^{a}\delta^2\frac{1}{2a}\dd\delta
=\frac{1}{2a}\left[\frac{\delta^3}{3}\right]_{-a}^{a}
=\frac{a^2}{3},\qquad
\sigma=\frac{a}{\sqrt{3}} .
\end{equation*}

\subsection{反正弦分布}
反正弦分布为
\begin{equation*}
f(\delta)=\frac{1}{\pi\sqrt{a^2-\delta^2}},\qquad \abs{\delta}<a .
\end{equation*}
该分布关于零点对称。令 $\delta=a\sin\theta$，则 $\dd\delta=a\cos\theta\dd\theta$，$\sqrt{a^2-\delta^2}=a\cos\theta$。当 $\delta$ 从 $-a$ 到 $a$ 时，$\theta$ 从 $-\pi/2$ 到 $\pi/2$。因此
\begin{equation*}
\sigma^2=\int_{-a}^{a}\delta^2\frac{1}{\pi\sqrt{a^2-\delta^2}}\dd\delta
=\frac{1}{\pi}\int_{-\pi/2}^{\pi/2}a^2\sin^2\theta\dd\theta
=\frac{a^2}{2},\qquad
\sigma=\frac{a}{\sqrt{2}} .
\end{equation*}

\subsection{三角分布}
三角分布为
\begin{equation*}
f(\delta)=
\begin{cases}
\dfrac{a+\delta}{a^2}, & -a\le \delta<0,\\[4pt]
\dfrac{a-\delta}{a^2}, & 0\le \delta\le a .
\end{cases}
\end{equation*}
它关于零点对称，因此均值为零。方差为
\begin{equation*}
\sigma^2=2\int_{0}^{a}\delta^2\frac{a-\delta}{a^2}\dd\delta
=\frac{2}{a^2}\int_{0}^{a}(a\delta^2-\delta^3)\dd\delta
=\frac{a^2}{6},\qquad
\sigma=\frac{a}{\sqrt{6}} .
\end{equation*}

\subsection{直角分布}
直角分布为
\begin{equation*}
f(\delta)=\frac{\delta+a}{2a^2},\qquad -a\le \delta\le a .
\end{equation*}
该分布不关于零点对称，必须先求均值：
\begin{equation*}
\mu=\int_{-a}^{a}\delta\frac{\delta+a}{2a^2}\dd\delta
=\frac{1}{2a^2}\int_{-a}^{a}(\delta^2+a\delta)\dd\delta
=\frac{a}{3} .
\end{equation*}
再求二阶原点矩：
\begin{equation*}
\E(\delta^2)=\int_{-a}^{a}\delta^2\frac{\delta+a}{2a^2}\dd\delta
=\frac{1}{2a^2}\int_{-a}^{a}(\delta^3+a\delta^2)\dd\delta
=\frac{a^2}{3} .
\end{equation*}
于是
\begin{equation*}
\sigma^2=\E(\delta^2)-\mu^2=\frac{a^2}{3}-\frac{a^2}{9}=\frac{2a^2}{9},
\qquad
\sigma=\frac{\sqrt{2}a}{3} .
\end{equation*}

\subsection{椭圆分布}
椭圆分布为
\begin{equation*}
f(\delta)=\frac{2}{\pi a^2}\sqrt{a^2-\delta^2},\qquad \abs{\delta}\le a .
\end{equation*}
它关于零点对称。令 $\delta=a\sin\theta$，则
\begin{equation*}
\sigma^2=\int_{-a}^{a}\delta^2\frac{2}{\pi a^2}\sqrt{a^2-\delta^2}\dd\delta
=\frac{2a^2}{\pi}\int_{-\pi/2}^{\pi/2}\sin^2\theta\cos^2\theta\dd\theta .
\end{equation*}
由 $\sin^2\theta\cos^2\theta=\dfrac{1}{8}(1-\cos4\theta)$，得
\begin{equation*}
\int_{-\pi/2}^{\pi/2}\sin^2\theta\cos^2\theta\dd\theta=\frac{\pi}{8}.
\end{equation*}
因此
\begin{equation*}
\sigma^2=\frac{2a^2}{\pi}\cdot\frac{\pi}{8}=\frac{a^2}{4},
\qquad
\sigma=\frac{a}{2} .
\end{equation*}

\subsection{双三角分布}
双三角分布为
\begin{equation*}
f(\delta)=\frac{\abs{\delta}}{a^2},\qquad \abs{\delta}\le a .
\end{equation*}
它关于零点对称。方差为
\begin{equation*}
\sigma^2=2\int_{0}^{a}\delta^2\frac{\delta}{a^2}\dd\delta
=\frac{2}{a^2}\int_{0}^{a}\delta^3\dd\delta
=\frac{a^2}{2},\qquad
\sigma=\frac{a}{\sqrt{2}} .
\end{equation*}

\newpage

\section{系统误差与粗大误差判别公式说明}

\subsection{两个平均值比较公式}
设两组测量平均值分别为 $\bar{x}_i$ 和 $\bar{x}_j$，其标准差分别为 $\sigma_i$ 和 $\sigma_j$，并假定二者独立。考虑差值
\begin{equation*}
    D=\bar{x}_i-\bar{x}_j .
\end{equation*}
若两组数据之间不存在显著系统差异，则 $D$ 的理论均值应为零。由独立变量方差相加得
\begin{equation*}
    \Var(D)=\Var(\bar{x}_i)+\Var(\bar{x}_j)=\sigma_i^2+\sigma_j^2 .
\end{equation*}
因此差值的标准差为
\begin{equation*}
    \sigma_D=\sqrt{\sigma_i^2+\sigma_j^2} .
\end{equation*}
若采用约 $2\sigma$ 的经验界限，则比较量为
\begin{equation*}
    \abs{\bar{x}_i-\bar{x}_j}
    \quad \text{与}\quad
    2\sqrt{\sigma_i^2+\sigma_j^2} .
\end{equation*}
若差值超过该界限，通常认为两组平均值之间存在显著差异；若没有超过，则认为差异可由随机误差解释。

\subsection{秩和检验大样本近似推导}
将两组数据合并后按大小排序，并给出秩次。设第一组样本量为 $n_1$，第二组样本量为 $n_2$，第一组秩和为 $T$。在无系统差异的原假设下，第一组相当于从 $1,2,\ldots,n_1+n_2$ 这些秩次中随机抽取 $n_1$ 个。

设总样本量
\begin{equation*}
    N=n_1+n_2 .
\end{equation*}
所有秩次的平均值为
\begin{equation*}
    \frac{1+N}{2} .
\end{equation*}
故秩和的数学期望为
\begin{equation*}
    a=\E(T)=n_1\frac{N+1}{2}
    =\frac{n_1(n_1+n_2+1)}{2} .
\end{equation*}
不放回抽样秩和的方差为
\begin{equation*}
    \sigma_T^2=\frac{n_1n_2(N+1)}{12}
    =\frac{n_1n_2(n_1+n_2+1)}{12} .
\end{equation*}
当 $n_1,n_2$ 较大时，秩和近似服从正态分布，因此标准化统计量为
\begin{equation*}
    Z=\frac{T-a}{\sigma_T}
    =\frac{T-\dfrac{n_1(n_1+n_2+1)}{2}}
    {\sqrt{\dfrac{n_1n_2(n_1+n_2+1)}{12}}} .
\end{equation*}
若
\begin{equation*}
    \abs{Z}>t_\alpha,
\end{equation*}
则认为两组数据的中心位置差异显著，从而可判断存在系统差异。

\newpage

\subsection{两样本 $t$ 检验公式推导}
设两组样本均值分别为 $\bar{x},\bar{y}$，样本容量分别为 $n_x,n_y$，样本标准差分别为 $\sigma_x,\sigma_y$。若认为两组总体方差相同，则合并方差估计通常为
\begin{equation*}
    s_p^2=\frac{n_x\sigma_x^2+n_y\sigma_y^2}{n_x+n_y-2} .
\end{equation*}
这里分母为 $n_x+n_y-2$，是因为两组数据各自估计了一个均值，共损失两个自由度。

均值差的标准差为
\begin{equation*}
    s_{\bar{x}-\bar{y}}=s_p\sqrt{\frac{1}{n_x}+\frac{1}{n_y}}
    =s_p\sqrt{\frac{n_x+n_y}{n_xn_y}} .
\end{equation*}
所以
\begin{equation*}
    \abs{t}
    =\frac{\abs{\bar{x}-\bar{y}}}{s_p\sqrt{(n_x+n_y)/(n_xn_y)}}
    =\abs{\bar{x}-\bar{y}}\sqrt{\frac{n_xn_y(n_x+n_y-2)}{(n_x+n_y)(n_x\sigma_x^2+n_y\sigma_y^2)}} .
\end{equation*}
若 $\abs{t}\ge t_\alpha$，则认为两组均值差异显著。

\subsection{莱以特准则的概率依据}
若误差服从正态分布，则
\begin{equation*}
    P(\abs{\delta}\le 3\sigma)\approx 0.9973 .
\end{equation*}
这说明正常随机误差落在 $[-3\sigma,3\sigma]$ 外的概率约为 $0.27\%$。所以若某个残余误差满足
\begin{equation*}
    \abs{v_i}>3\sigma,
\end{equation*}
则它很难仅由正常随机波动产生，因而可怀疑为粗大误差。实际使用时，$\sigma$ 常用贝塞尔公式估计：
\begin{equation*}
    \sigma=\sqrt{\frac{\sum_{i=1}^{n}v_i^2}{n-1}} .
\end{equation*}

\subsection{罗曼诺夫斯基准则的留一思想}
判断第 $j$ 个观测值 $x_j$ 是否异常时，为避免异常值本身影响平均值，常把它从样本中临时剔除，用其余 $n-1$ 个数据计算平均值：
\begin{equation*}
    \bar{x}_{(-j)}=\frac{1}{n-1}\sum_{i=1,i\ne j}^{n}x_i .
\end{equation*}
再用其余数据估计标准差。因为剔除了一个可疑值，还剩 $n-1$ 个数据，并估计了一个平均值，所以自由度为 $n-2$：
\begin{equation*}
    s_{(-j)}=\sqrt{\frac{\sum_{i=1}^{n}v_i^2}{n-2}} .
\end{equation*}
于是构造标准化偏离量
\begin{equation*}
    \frac{\abs{x_j-\bar{x}_{(-j)}}}{s_{(-j)}}
    =
    \frac{
    \left|x_j-\dfrac{1}{n-1}\sum_{i=1,i\ne j}^{n}x_i\right|
    }{
    \sqrt{\dfrac{\sum_{i=1}^{n}v_i^2}{n-2}}
    } .
\end{equation*}
若该值超过临界值 $K(n,\alpha)$，则判为粗大误差。

\newpage

\subsection{格罗布斯准则与狄克松准则的统计量含义}
格罗布斯准则先将数据从小到大排列为
\begin{equation*}
    x_{(1)}\le x_{(2)}\le\cdots\le x_{(n)} .
\end{equation*}
若怀疑最小值异常，则考察
\begin{equation*}
    g_{(1)}=\frac{\bar{x}-x_{(1)}}{\sigma};
\end{equation*}
若怀疑最大值异常，则考察
\begin{equation*}
    g_{(n)}=\frac{x_{(n)}-\bar{x}}{\sigma} .
\end{equation*}
其本质是把端点偏离平均值的距离用标准差标准化，再与临界值 $g_0(n,\alpha)$ 比较。

狄克松准则不计算标准差，而是用“可疑端点与邻近点的间隔”除以“样本总体跨度”或修正后的跨度。例如当 $n\le 7$ 时，最小端异常的统计量为
\begin{equation*}
    r_{10}=\frac{x_{(2)}-x_{(1)}}{x_{(n)}-x_{(1)}} .
\end{equation*}
若 $x_{(1)}$ 远离其他数据，则 $x_{(2)}-x_{(1)}$ 大，$r_{10}$ 也大。最大端异常时可写为
\begin{equation*}
    r_{10}=\frac{x_{(n)}-x_{(n-1)}}{x_{(n)}-x_{(1)}} .
\end{equation*}
样本量变大后，分母需要排除一个或两个端点以降低另一端异常值的干扰，因此出现 \\$r_{11},r_{21},r_{22}$ 等不同形式。

\section{误差合成与误差分配推导}

\subsection{函数标准差合成公式推导}
设输入误差为随机量 $\Delta x_i$，输出误差的一阶近似为
\begin{equation*}
    \Delta y=\sum_{i=1}^{n}a_i\Delta x_i,
    \qquad
    a_i=\frac{\partial f}{\partial x_i} .
\end{equation*}
输出方差为
\begin{equation*}
    \sigma_y^2=\Var(\Delta y)=\Var\left(\sum_{i=1}^{n}a_i\Delta x_i\right) .
\end{equation*}
展开方差：
\begin{equation*}
    \Var\left(\sum_{i=1}^{n}a_i\Delta x_i\right)
    =\sum_{i=1}^{n}a_i^2\Var(\Delta x_i)
    +2\sum_{1\le i<j\le n}a_i a_j\Cov(\Delta x_i,\Delta x_j) .
\end{equation*}
令
\begin{equation*}
    \Var(\Delta x_i)=\sigma_{x_i}^2,
    \qquad
    \Cov(\Delta x_i,\Delta x_j)=\rho_{ij}\sigma_{x_i}\sigma_{x_j},
\end{equation*}
其中 $\rho_{ij}$ 为相关系数。代入得
\begin{equation*}
    \sigma_y^2
    =\sum_{i=1}^{n}a_i^2\sigma_{x_i}^2
    +2\sum_{1\le i<j\le n}a_i a_j\rho_{ij}\sigma_{x_i}\sigma_{x_j} .
\end{equation*}
再开方：
\begin{equation*}
    \sigma_y=\bigg[
    \sum_{i=1}^{n}\left(\frac{\partial f}{\partial x_i}\right)^2\sigma_{x_i}^2
    +2\sum_{1\le i<j\le n}
    \frac{\partial f}{\partial x_i}\frac{\partial f}{\partial x_j}
    \rho_{ij}\sigma_{x_i}\sigma_{x_j}
    \bigg]^{1/2} .
\end{equation*}

\subsection{等作用原则误差分配推导}
设输出函数一阶误差为
\begin{equation*}
    \Delta y=\sum_{i=1}^{n}a_i\Delta x_i .
\end{equation*}
若各输入互不相关，则
\begin{equation*}
    \sigma_y^2=\sum_{i=1}^{n}a_i^2\sigma_i^2 .
\end{equation*}
等作用原则要求每个误差源对输出方差的贡献相等，即
\begin{equation*}
    a_i^2\sigma_i^2=\frac{\sigma_y^2}{n} .
\end{equation*}
两边开方，得
\begin{equation*}
    \abs{a_i}\sigma_i=\frac{\sigma_y}{\sqrt{n}} .
\end{equation*}
若取 $a_i>0$ 或用绝对值表示灵敏度，则
\begin{equation*}
    \sigma_i=\frac{\sigma_y}{\sqrt{n}}\frac{1}{\abs{a_i}}
    =\frac{\sigma_y}{\sqrt{n}}\frac{1}{\abs{\partial f/\partial x_i}} .
\end{equation*}
同理，对极限误差有
\begin{equation*}
    \delta_i=\frac{\delta}{\sqrt{n}}\frac{1}{\abs{a_i}} .
\end{equation*}

\section{不确定度合成公式推导}

设测量结果由
\begin{equation*}
    y=f(x_1,x_2,\ldots,x_N)
\end{equation*}
给出。对函数作一阶展开：
\begin{equation*}
    \Delta y\approx \sum_{i=1}^{N}\frac{\partial f}{\partial x_i}\Delta x_i .
\end{equation*}
令输入量的标准不确定度为 $u_{x_i}$，灵敏系数为
\begin{equation*}
    c_i=\frac{\partial f}{\partial x_i} .
\end{equation*}
若各输入量独立，则
\begin{equation*}
    u_c^2(y)=\sum_{i=1}^{N}c_i^2u_{x_i}^2 .
\end{equation*}
记每个输入量对输出量贡献的标准不确定度为
\begin{equation*}
    u_i=c_i u_{x_i},
\end{equation*}
则
\begin{equation*}
    u_c=\sqrt{\sum_{i=1}^{N}\left(\frac{\partial f}{\partial x_i}\right)^2u_{x_i}^2}
    =\sqrt{\sum_{i=1}^{N}u_i^2} .
\end{equation*}
这与误差合成公式形式完全一致，只是把标准差符号换成标准不确定度符号。

\newpage

\section{最小二乘法公式推导}

\subsection{线性误差方程}
设有 $n$ 个观测方程，$t$ 个未知参数。写成矩阵形式为
\begin{equation*}
    L=A X+\varepsilon,
\end{equation*}

\begin{equation*}
    L=\begin{pmatrix}l_1\\l_2\\\vdots\\l_n\end{pmatrix},
    \qquad
    X=\begin{pmatrix}x_1\\x_2\\\vdots\\x_t\end{pmatrix},
    \qquad
    A=(a_{ij})_{n\times t} .
\end{equation*}
用估计量 $\hat{X}$ 代替真参数后，拟合值为 $A\hat{X}$，残差向量为 $V=L-A\hat{X}$。

\subsection{等精度最小二乘正规方程}
等精度时，各观测方程可靠程度相同，应使残差平方和最小：
\begin{equation*}
    \Phi=V^\T V=(L-A\hat{X})^\T(L-A\hat{X}) .
\end{equation*}
展开：
\begin{equation*}
\begin{aligned}
    \Phi
    &=L^\T L-L^\T A\hat{X}-\hat{X}^\T A^\T L+
      \hat{X}^\T A^\T A\hat{X} .
\end{aligned}
\end{equation*}
由于 $L^\T A\hat{X}$ 是标量，等于其转置 $\hat{X}^\T A^\T L$，所以
\begin{equation*}
    \Phi=L^\T L-2\hat{X}^\T A^\T L+
    \hat{X}^\T A^\T A\hat{X} .
\end{equation*}
对 $\hat{X}$ 求导并令其为零：
\begin{equation*}
    \frac{\partial\Phi}{\partial \hat{X}}
    =-2A^\T L+2A^\T A\hat{X}=0 .
\end{equation*}
故
\begin{equation*}
    A^\T A\hat{X}=A^\T L .
\end{equation*}
这就是正规方程。若 $A^\T A$ 可逆，则
\begin{equation*}
    \hat{X}=(A^\T A)^{-1}A^\T L .
\end{equation*}
记 $C=A^\T A$，则 $\hat{X}=C^{-1}A^\T L$。

\subsection{不等精度最小二乘正规方程}
不等精度时，各观测值的方差不同。方差小的观测值应有更大权重。设权阵为
\begin{equation*}
    P=\diag(p_1,p_2,\ldots,p_n)=
    \begin{bmatrix}
p_1 & 0 & \cdots & 0\\
0 & p_2 & \cdots & 0\\
\vdots & \vdots & \ddots & \vdots\\
0 & 0 & \cdots & p_n
\end{bmatrix}.
\end{equation*}
此时应最小化加权残差平方和
\begin{equation*}
    \Phi=V^\T P V=(L-A\hat{X})^\T P(L-A\hat{X}) .
\end{equation*}
展开并求导：
\begin{equation*}
\begin{aligned}
    \Phi
    &=L^\T P L-2\hat{X}^\T A^\T P L+
      \hat{X}^\T A^\T P A\hat{X},\\
    \frac{\partial\Phi}{\partial\hat{X}}
    &=-2A^\T P L+2A^\T P A\hat{X}=0 .
\end{aligned}
\end{equation*}
于是
\begin{equation*}
    A^\T P A\hat{X}=A^\T P L .
\end{equation*}
若 $A^\T P A$ 可逆，则
\begin{equation*}
    \hat{X}=(A^\T P A)^{-1}A^\T P L .
\end{equation*}
记 $C=A^\T P A$，则 $\hat{X}=C^{-1}A^\T P L$。等精度是 $P=I$ 的特殊情况。

\subsection{单位权标准差与自由度}
设共有 $n$ 个观测方程、$t$ 个未知参数。最小二乘中，$t$ 个参数由数据估计出来，因此残差的独立个数为
\begin{equation*}
    r=n-t .
\end{equation*}
这就是自由度。等精度时，单位权标准差估计为
\begin{equation*}
    \sigma_0=\sqrt{\frac{V^\T V}{n-t}}
    =\sqrt{\frac{\sum_{i=1}^{n}v_i^2}{n-t}} .
\end{equation*}
不等精度时，采用加权残差平方和：
\begin{equation*}
    \sigma_0=\sqrt{\frac{V^\T P V}{n-t}}
    =\sqrt{\frac{\sum_{i=1}^{n}p_i v_i^2}{n-t}} .
\end{equation*}

\subsection{参数估计量标准差推导}
在不等精度线性模型中，有
\begin{equation*}
    \hat{X}=(A^\T P A)^{-1}A^\T P L .
\end{equation*}
若观测误差协方差阵与权阵满足
\begin{equation*}
    \Cov(L)=\sigma_0^2P^{-1},
\end{equation*}
则估计量协方差阵为
\begin{equation*}
\begin{aligned}
    \Cov(\hat{X})
    &=(A^\T P A)^{-1}A^\T P\Cov(L)P A(A^\T P A)^{-1} \\
    &=(A^\T P A)^{-1}A^\T P(\sigma_0^2P^{-1})P A(A^\T P A)^{-1} \\
    &=\sigma_0^2(A^\T P A)^{-1} .
\end{aligned}
\end{equation*}
设
\begin{equation*}
    (A^\T P A)^{-1}=
    \begin{pmatrix}
    d_{11}&d_{12}&\cdots&d_{1t}\\
    d_{21}&d_{22}&\cdots&d_{2t}\\
    \vdots&\vdots&\ddots&\vdots\\
    d_{t1}&d_{t2}&\cdots&d_{tt}
    \end{pmatrix} .
\end{equation*}
则第 $i$ 个估计量的方差为
\begin{equation*}
    \sigma_{x_i}^2=\sigma_0^2 d_{ii} .
\end{equation*}
所以
\begin{equation*}
    \sigma_{x_i}=\sigma_0\sqrt{d_{ii}},
    \qquad i=1,2,\ldots,t .
\end{equation*}

\newpage

\section{一元线性回归公式推导}

\subsection{最小二乘目标函数}
设一元线性回归模型为
\begin{equation*}
    \hat{y}_t=b_0+bx_t,
    \qquad t=1,2,\ldots,N .
\end{equation*}
残差为
\begin{equation*}
    e_t=y_t-\hat{y}_t=y_t-b_0-bx_t .
\end{equation*}
最小二乘法要求残差平方和最小：
\begin{equation*}
    Q(b_0,b)=\sum_{t=1}^{N}(y_t-b_0-bx_t)^2 .
\end{equation*}

\subsection{回归系数与截距推导}
对 $b_0$ 和 $b$ 分别求偏导：
\begin{equation*}
\begin{aligned}
    \frac{\partial Q}{\partial b_0}
    &=-2\sum_{t=1}^{N}(y_t-b_0-bx_t)=0,\\
    \frac{\partial Q}{\partial b}
    &=-2\sum_{t=1}^{N}x_t(y_t-b_0-bx_t)=0 .
\end{aligned}
\end{equation*}
化简第一式：
\begin{equation*}
    \sum_{t=1}^{N}y_t-Nb_0-b\sum_{t=1}^{N}x_t=0 .
\end{equation*}
两边除以 $N$，得
\begin{equation*}
    \bar{y}-b_0-b\bar{x}=0,
    \qquad
    b_0=\bar{y}-b\bar{x} .
\end{equation*}
代入第二个正规方程：
\begin{equation*}
    \sum_{t=1}^{N}x_t\left[y_t-(\bar{y}-b\bar{x})-bx_t\right]=0 .
\end{equation*}
即
\begin{equation*}
    \sum_{t=1}^{N}x_t(y_t-\bar{y})-b\sum_{t=1}^{N}x_t(x_t-\bar{x})=0 .
\end{equation*}
注意
\begin{equation*}
    \sum_{t=1}^{N}x_t(y_t-\bar{y})
    =\sum_{t=1}^{N}(x_t-\bar{x})(y_t-\bar{y})=l_{xy},
\end{equation*}
以及
\begin{equation*}
    \sum_{t=1}^{N}x_t(x_t-\bar{x})
    =\sum_{t=1}^{N}(x_t-\bar{x})^2=l_{xx} .
\end{equation*}
所以
\begin{equation*}
    l_{xy}-b l_{xx}=0,
    \qquad
    b=\frac{l_{xy}}{l_{xx}} .
\end{equation*}
再代回得到
\begin{equation*}
    b_0=\bar{y}-b\bar{x} .
\end{equation*}

\newpage

\subsection{离差平方和恒等式}
自变量离差平方和为
\begin{equation*}
    l_{xx}=\sum_{t=1}^{N}(x_t-\bar{x})^2 .
\end{equation*}
展开：
\begin{equation*}
\begin{aligned}
    l_{xx}
    &=\sum_{t=1}^{N}(x_t^2-2x_t\bar{x}+\bar{x}^2) \\
    &=\sum_{t=1}^{N}x_t^2-2\bar{x}\sum_{t=1}^{N}x_t+N\bar{x}^2 .
\end{aligned}
\end{equation*}
由于 $\sum x_t=N\bar{x}$，所以
\begin{equation*}
    l_{xx}=\sum_{t=1}^{N}x_t^2-2N\bar{x}^2+N\bar{x}^2
    =\sum_{t=1}^{N}x_t^2-N\bar{x}^2 .
\end{equation*}
又
\begin{equation*}
    N\bar{x}^2=N\left(\frac{1}{N}\sum_{t=1}^{N}x_t\right)^2
    =\frac{1}{N}\left(\sum_{t=1}^{N}x_t\right)^2 .
\end{equation*}
故
\begin{equation*}
    l_{xx}=\sum_{t=1}^{N}x_t^2-\frac{1}{N}\left(\sum_{t=1}^{N}x_t\right)^2 .
\end{equation*}
同理可得
\begin{equation*}
    l_{yy}=\sum_{t=1}^{N}y_t^2-\frac{1}{N}\left(\sum_{t=1}^{N}y_t\right)^2 .
\end{equation*}

\subsection{离差积和恒等式}
离差积和为
\begin{equation*}
    l_{xy}=\sum_{t=1}^{N}(x_t-\bar{x})(y_t-\bar{y}) .
\end{equation*}
展开：
\begin{equation*}
\begin{aligned}
    l_{xy}
    &=\sum_{t=1}^{N}\left(x_ty_t-x_t\bar{y}-\bar{x}y_t+\bar{x}\bar{y}\right) \\
    &=\sum_{t=1}^{N}x_ty_t-\bar{y}\sum_{t=1}^{N}x_t-\bar{x}\sum_{t=1}^{N}y_t+N\bar{x}\bar{y} .
\end{aligned}
\end{equation*}
由 $\sum x_t=N\bar{x}$、$\sum y_t=N\bar{y}$，得
\begin{equation*}
    l_{xy}=\sum_{t=1}^{N}x_ty_t-N\bar{x}\bar{y} .
\end{equation*}
而
\begin{equation*}
    N\bar{x}\bar{y}=\frac{1}{N}\left(\sum_{t=1}^{N}x_t\right)
    \left(\sum_{t=1}^{N}y_t\right) .
\end{equation*}
所以
\begin{equation*}
    l_{xy}=\sum_{t=1}^{N}x_ty_t-\frac{1}{N}\left(\sum_{t=1}^{N}x_t\right)
    \left(\sum_{t=1}^{N}y_t\right) .
\end{equation*}

\subsection{平方和分解推导}
总离差平方和为
\begin{equation*}
    S=\sum_{t=1}^{N}(y_t-\bar{y})^2=l_{yy} .
\end{equation*}
将 $y_t-\bar{y}$ 分解为
\begin{equation*}
    y_t-\bar{y}=(\hat{y}_t-\bar{y})+(y_t-\hat{y}_t) .
\end{equation*}
于是
\begin{equation*}
\begin{aligned}
    S
    &=\sum_{t=1}^{N}\left[(\hat{y}_t-\bar{y})+(y_t-\hat{y}_t)\right]^2 \\
    &=\sum_{t=1}^{N}(\hat{y}_t-\bar{y})^2
      +\sum_{t=1}^{N}(y_t-\hat{y}_t)^2
      +2\sum_{t=1}^{N}(\hat{y}_t-\bar{y})(y_t-\hat{y}_t) .
\end{aligned}
\end{equation*}
最小二乘回归中，残差与拟合值正交，因此
\begin{equation*}
    \sum_{t=1}^{N}(\hat{y}_t-\bar{y})(y_t-\hat{y}_t)=0 .
\end{equation*}
所以
\begin{equation*}
    S=U+Q,
\end{equation*}
其中
\begin{equation*}
    U=\sum_{t=1}^{N}(\hat{y}_t-\bar{y})^2,
    \qquad
    Q=\sum_{t=1}^{N}(y_t-\hat{y}_t)^2 .
\end{equation*}
由于
\begin{equation*}
    \hat{y}_t-\bar{y}=b_0+bx_t-\bar{y}=b(x_t-\bar{x}),
\end{equation*}
因此
\begin{equation*}
    U=\sum_{t=1}^{N}b^2(x_t-\bar{x})^2=b^2l_{xx} .
\end{equation*}
又 $b=l_{xy}/l_{xx}$，故
\begin{equation*}
    U=b^2l_{xx}=b\,l_{xy} .
\end{equation*}
因此
\begin{equation*}
    Q=S-U=l_{yy}-b\,l_{xy} .
\end{equation*}

\subsection{回归方差分析与 $F$ 统计量}
一元线性回归中，回归平方和对应一个回归系数 $b$，故自由度为 $\nu_U=1$。总平方和自由度为 $N-1$，因为估计了一个总体均值 $\bar{y}$。残差平方和在估计 $b_0,b$ 两个参数后还剩 $\nu_Q=N-2$ 个自由度。于是回归均方和残差均方分别为
\begin{equation*}
    M_U=\frac{U}{\nu_U}=\frac{U}{1},
    \qquad
    M_Q=\frac{Q}{\nu_Q}=\frac{Q}{N-2} .
\end{equation*}

\end{document}